\section{Introduction}
\label{sec:intro}

Driving benchmarks are built from clear, well-lit imagery. A policy that scores well on them
has been asked, almost exclusively, to act on scenes whose appearance is easy. What its
internal representation does when only appearance changes---and whether that change reaches
the planned trajectory---is not something a leaderboard number reports.

We study this directly. Starting from mined human braking events, we extract three directions
in each policy's intermediate representation: the \emph{hazard-response} direction
$v_{\mathrm{faith}}$, obtained from paired frames that differ only in whether the hazard is
present; the \emph{braking} direction $v_{\mathrm{decel}}$, obtained by contrasting cruising
and braking populations at matched speed; and the \emph{appearance} direction
$v_{\mathrm{bright}}$, obtained by re-rendering the same frame as night. We then ask how these
sit relative to one another, and what happens to the planned trajectory.

\paragraph{Finding.} The three directions do not spread out evenly. $v_{\mathrm{faith}}$ and
$v_{\mathrm{decel}}$ cluster ($34$--$46^\circ$; \Cref{tab:geometry}); $v_{\mathrm{bright}}$
sits on the opposite side of the random baseline from both. Because
$v_{\mathrm{faith}} = h(\text{hazard present}) - h(\text{hazard removed})$, a negative
projection means the model reads ``it got dark'' the way it reads ``the hazard went away.''
The behavioural readout agrees: synthetic night raises planned speed by $2.05$\,m/s in
DiffusionDrive, where deleting the lead vehicle from the image raises it by $0.23$\,m/s. Two
independent readouts---representation geometry and output trajectory---point to the same
mechanism. The shift is concentrated in two of four policies ($|\mu|/\sigma = 2.12$ and $1.34$,
sign agreement $>94\%$) and dispersed in the other two, where the distribution is bimodal and
the mean describes no event; we report both cases rather than averaging over the difference.

\paragraph{What carries it.} Darkness alone is nearly inert; the effect is carried by sensor
noise. Perturbing only the sky, with the driving corridor untouched, still recovers a fifth of
the hazard-removal effect. The dependence is therefore global, not mediated by a local cue in
the region that matters for driving.

\paragraph{What we could not show.} A natural next step is to use the projection onto these
directions as a cheap behavioural proxy---to rank policies, or to flag individual events,
without closed-loop simulation. We pre-registered five such tests. Four are rejected, and the
fifth holds for only two of four policies (\Cref{tab:negative}). We report this in full: the
practice of treating a representation-space projection as a behavioural proxy is common, and
our evidence says it does not hold up here. One readout does transfer across the domain
shift---a multi-layer speed axis (\Cref{tab:speed})---and it is instructive that the one that
works has an \emph{external} target variable.

\paragraph{Contributions.}
\begin{enumerate}\itemsep2pt
\item An appearance-induced failure mode, consistent across four end-to-end policies, in which
      a night rendering moves the representation opposite to the policy's own hazard direction
      and raises planned speed (\Cref{sec:geometry,sec:behaviour}).
\item Isolation of its cause: sensor noise rather than luminance, and a global rather than
      local dependence---established by controls that hold the driving corridor fixed
      (\Cref{sec:controls}).
\item Five pre-registered tests of direction projections as behavioural proxies, four
      rejected, with the discipline that made the rejections interpretable
      (\Cref{sec:negative}).
\end{enumerate}
